\subsection{トピックと参考文献の対応}

この表は、SWEBOK 16章が示すトピック対参考資料の行列を、学習用に簡略化したものである。詳しく学ぶ場合は、各トピックに対応する文献を読むとよい。

\par\smallskip\noindent\refstepcounter{table}\label{tab:ch16-28}表 \thetable: トピックと参考文献の対応におけるトピック・主な参考資料\par\smallskip
表 \thetable は、トピックと参考文献の対応におけるトピック・主な参考資料を比較・確認しやすい形で整理したものである。\par\smallskip
\begin{longtable}{>{\raggedright\arraybackslash}p{0.36\textwidth}>{\raggedright\arraybackslash}p{0.58\textwidth}}
\toprule
トピック & 主な参考資料 \\
\midrule
\endfirsthead
\toprule
トピック & 主な参考資料 \\
\midrule
\endhead
1 システムまたは解決策の基本概念 & Sommerville, Software Engineering。 \\
2 コンピュータアーキテクチャと組織 & Null and Lobur, The Essentials of Computer Organization and Architecture; Sommerville。 \\
3 データ構造とアルゴリズム & Cormen et al., Introduction to Algorithms; Horowitz et al., Computer Algorithms; Null and Lobur。 \\
4 プログラミング基礎と言語 & Brookshear, Computer Science: An Overview; Null and Lobur; McConnell, Code Complete。 \\
5 基本ソフトウェア & Tanenbaum and Bos, Modern Operating Systems; Brookshear。 \\
6 データベース管理 & Connolly and Begg, Database Systems; Hernandez, Database Design for Mere Mortals; Fagin。 \\
7 コンピュータネットワークと通信 & Kurose and Ross, Computer Networking: A Top-Down Approach; Brookshear。 \\
8 人間要因 & Nielsen, Usability Engineering; McConnell; ISO 9241-420。 \\
9 人工知能と機械学習 & Goodfellow et al., Deep Learning; Shafiq et al.; Martínez-Fernández et al.; Washizaki et al.。 \\
\bottomrule
\end{longtable}
